% SGA 6, Exposición X: fragmento español alineado con la fuente.
% Alcance: sección 2 completa (Polinomios de Snapper).
% Fuente de trabajo francesa: sga6_fr_r1/src/fr.tex, líneas 12688--12887.
% Testigo impreso: Théorie des Intersections et Théorème de Riemann--Roch,
% Exposé X, pp. 531--538. Las lecturas diplomáticas dudosas se conservan
% como primer argumento de \sisi y la lectura matemática corregida como segundo.

\providecommand{\sisi}[2]{#2}

\section*{2. Polinomios de Snapper}

\subsection*{2.1}
Recordemos primero la definición y algunas propiedades elementales de las
aplicaciones polinómicas:

\textbf{Definición 2.1.1.} Sea $f$ una aplicación de un grupo abeliano $A$
en un grupo abeliano $B$. Se dice que $f$ es una aplicación polinómica si,
para toda sucesión $x_1,\ldots,x_r$ ($r\geq1$) de elementos de $A$, existe
una familia
$(y(x_1,\ldots,x_r;k_1,\ldots,k_r))_{k_1,\ldots,k_r\geq0}$ de elementos de
$B$, casi todos nulos, que satisface la condición siguiente: para toda
sucesión $n_1,\ldots,n_r$ de enteros, se tiene
\[
f(n_1x_1+\cdots+n_rx_r)=
\sum_{k_1,\ldots,k_r\geq0}(n_1^{k_1}\cdots n_r^{k_r})
y(x_1,\ldots,x_r;k_1,\ldots,k_r).
\]
Se dice que $f$ es de grado $\leq m$ si, cualesquiera que sean
$r\in\sisi{\mathbb Z}{\mathbb Z_{\geq1}}$ y
$x_1,\ldots,x_r\in A$, se tiene
\[
y(x_1,\ldots,x_r;k_1,\ldots,k_r)=0
\]
para $k_1+\cdots+k_r>m$.

\textbf{Lema 2.1.2.} Todo homomorfismo de grupos abelianos es una aplicación
polinómica de grado $\leq1$. Si $h=f\circ g$, donde $f$ es una aplicación
polinómica de grado $p$ y $g$ una aplicación polinómica de grado $q$, entonces
$h$ es una aplicación polinómica de grado $\leq pq$.

Demostración trivial.

\textbf{Lema 2.1.3.} Sea
\[
0\longrightarrow G'\longrightarrow G\xrightarrow{j}G''\longrightarrow0
\]
una sucesión exacta de grupos abelianos tal que $G'$ es divisible por un
número primo $p$ y $G''$ es un grupo de torsión. Entonces toda aplicación
polinómica $f:G\to\mathbb Q$ que se factoriza a través de $\mathbb Z$ es
constante.

\emph{Demostración:} Se trata de demostrar que, para todo $x\in G$, se tiene
$f(x)=f(0)$.

a) Supongamos primero que $x$ es un elemento de torsión, digamos $mx=0$
($m\in Z$, $m\ne0$). Entonces el polinomio
$P_x(n)=f(nx)-f(0)$ tiene infinitas raíces, a saber, todos los enteros
divisibles por $m$. Por tanto, el polinomio $P_x$ es idénticamente nulo y se
tiene $f(x)=f(0)$.

b) Supongamos ahora que $x$ es un elemento libre de $G$ divisible por todas
las potencias $p^k$ ($k\in\mathbb Z$). Consideremos los polinomios $P_k$
definidos por
\[
P_k(n)=f(nx/p^k)=\sum_{i\in\mathbb N}n^iy(x/p^k;i)\in\mathbb Z
\qquad(n\in\mathbb Z;\ y(x/p^k,i)\in\mathbb Q).
\]
Para todo $n\in\mathbb Z$ se tiene la identidad
\[
P_k(n)=P_{k+1}(pn),
\]
es decir,
\[
\sum_{i\in\mathbb N}n^iy(x/p^k;i)=
\sum_{i\in\mathbb N}p^in^iy(x/p^{k+1};i),
\]
de donde, comparando coeficientes, se obtiene
\[
y(x/p^k;i)=p^iy(x/p^{k+1};i)\qquad(i\in\mathbb N).
\]
Ahora bien, todo polinomio $P:Q\to\mathbb Q$ es continuo para la
topología inducida por $\mathbb R$. Por tanto, para $k$ suficientemente
grande se tiene
\[
|P_x(n/p^k)-P_x(0)|<1.
\]
Pero, por definición, $P_x(n/p^k)$ es un entero para todos
$n,k\in\mathbb Z$. Se sigue que $P_x(n/p^k)-P_x(0)=0$ para $k$
suficientemente grande. Así pues, el polinomio $P_x$ es constante y, en
particular, $P_x(1)=P_x(0)$, es decir, $f(x)=f(0)$.

Los dos casos a) y b) anteriores muestran que la restricción de $f$ a $G'$
es constante. Mostremos ahora que $f$ es constante en cada clase de $G$
módulo $G'$. Para todo $b\in G$, denotemos por
$g_b:G'\to\underline Q$ la aplicación definida por
\[
g_b(y)=f(b+y)\qquad(y\in G').
\]
La aplicación $g_b$ es una aplicación polinómica de $G'$ en $\mathbb Q$ que
se factoriza a través de $\mathbb Z$ y, por lo visto anteriormente, es
constante. De aquí se sigue que $f$ es constante en las clases de $G$ módulo
$\sisi{G''}{G'}$. Por consiguiente, se puede escribir $f=h\circ j$, donde
$\sisi{h:G'\to\mathbb Q}{h:G''\to\mathbb Q}$ es una aplicación polinómica que se factoriza a través
de $\mathbb Z$. Pero, por hipótesis, $\sisi{C''}{G''}$ es un grupo de
torsión, de modo que $h$ es constante. Por tanto, $f$ es constante,
C.Q.F.D.

\textbf{Lema 2.1.4.} Sea $A$ un anillo conmutativo unitario e $I$ un ideal
nilpotente de $A$. Sea $P$ el subgrupo multiplicativo de $A^*$ formado por
los elementos invertibles de $A$ de la forma $x=1+y$, con $y\in I$, y sea
$i:P\to A$ la inclusión natural. Entonces la aplicación compuesta
\[
\alpha:P\xrightarrow{i}A\xrightarrow{j}A\otimes_{\mathbb Z}\mathbb Q,
\]
(donde $j(x)=x\otimes1$) es una aplicación polinómica. Si $I^{m+1}=0$,
entonces $\alpha$ es de grado $\leq m$.

\emph{Demostración:} Sea $x_1,\ldots,x_r$ una sucesión de elementos de $P$
con $x_i=1+y_i$ ($i=1,\ldots,r$). Entonces la fórmula binomial da, para
todos $n_1,\ldots,n_r\in Z$:
\begin{equation}\tag{2.1.5}
x_1^{n_1}\cdots x_r^{n_r}=
\sum_{k_1,\ldots,k_r\geq0}\binom{n_1}{k_1}\cdots\binom{n_r}{k_r}
y_1^{k_1}\cdots y_r^{k_r},
\end{equation}
donde la suma de la derecha es finita, pues los elementos
$y_1,\ldots,y_r$ son nilpotentes. Pero los coeficientes
$\binom{n_1}{k_1}\cdots\binom{n_r}{k_r}$ son polinomios con coeficientes
racionales en las variables $n_1,\ldots,n_r$, de grado
$k_1+\cdots+k_r$, por lo que $\alpha$ es efectivamente una aplicación
polinómica. Si $I^{m+1}=0$, entonces
$y_1^{k_1}\cdots y_r^{k_r}=0$ para $k_1+\cdots+k_r>m$. Por tanto,
$\alpha$ es de grado $\leq m$, C.Q.F.D.

\subsection*{2.2}
Consideremos ahora un esquema noetheriano $X$ y un elemento $y$ de
$\Filt_j(X)$ (1.1.1). Al asociar a todo $x\in K^\circ(X)$ el elemento $xy$
de $K_\bullet(X)$, se obtiene un homomorfismo de $K^\circ(X)$-módulos
\[
\Theta_y:K^\circ(X)\longrightarrow K_\bullet(X)
\]
que, teniendo en cuenta (1.3.2), se factoriza a través del epimorfismo
canónico
\[
j:K^\circ(X)\longrightarrow K^\circ(X)/\Filt^{j+1}(X)=A.
\]
Sea $\beta_y:A\to K_\bullet(X)$ el homomorfismo así definido. Si se denota
por $I$ el ideal $\sisi{\Filt^j(X)}{\Filt^1(X)}/\Filt^{j+1}(X)$ de $A$, se tiene $I^{j+1}=0$.
Conservando las notaciones del lema 2.1.4, se obtiene un homomorfismo de
grupos abelianos
\[
\gamma:\Pic(X)\longrightarrow P
\]
que asocia a la clase de todo $\mathcal O_X$-módulo invertible $L$ en
$\Pic(X)$ el elemento $j(\cl^\circ(L))$ de $P$.

De los lemas 2.1.2 y 2.1.4 se deduce inmediatamente:

\textbf{Proposición 2.2.1.} Para todo esquema noetheriano $X$ y para todo
$y\in\Filt_j(X)$ ($j\in\mathbb Z$), la aplicación compuesta
\[
\Pic(X)\longrightarrow K_\bullet(X)\otimes_{\mathbb Z}\mathbb Q,
\]
definida por
\[
\cl_P(L)\rightsquigarrow (\cl^\circ(L)y)\otimes1
\]
es una aplicación polinómica de grado $\leq j$.

\subsection*{2.3}
Si $X$ es un esquema propio sobre un cuerpo $k$ y
$g:X\to\Spec(k)$ es el morfismo estructural, entonces el homomorfismo
\[
g_*:K_\bullet(X)\longrightarrow K_\bullet(\Spec(k))\cong\mathbb Z
\]
asocia a la clase $\cl_\bullet(F)$ de todo $\mathcal O_X$-módulo coherente
$F$ la característica de Euler--Poincaré $\chi(F)\in\mathbb Z$ de $F$
sobre $X$. De 2.2.1 se deduce inmediatamente:

\textbf{Proposición 2.3.1.} Para todo esquema propio $X$ sobre $k$ y para
todo $y\in\Filt_j(X)$, la aplicación
\[
S_y:\Pic(X)\longrightarrow\mathbb Z\hookrightarrow\mathbb Q
\]
definida por
\[
\cl_P(L)\rightsquigarrow g_*(\cl^\circ(L)y)=\chi(\cl^\circ(L)y)
\]
es una aplicación polinómica de grado $\leq j$.

Cuando $L_1,\ldots,L_n$ son $\mathcal O_X$-módulos invertibles y $F$ es un
$\mathcal O_X$-módulo coherente tal que $\dim(\Supp(F))=j$, 2.3.1 implica,
en particular, que la característica de Euler--Poincaré
\begin{equation}\tag{2.3.2}
\chi(L_1^{\otimes m_1}\otimes\cdots\otimes L_n^{\otimes m_n}\otimes F)
\end{equation}
es un polinomio numérico en las variables $m_1,\ldots,m_n\in\mathbb Z$, de
grado $\leq\dimdot\operatorname{supp.}F$, denominado a veces polinomio de
Snapper para $F$ y la sucesión $L_1,\ldots,L_n$.

\subsection*{2.4}
Recordemos [2] que el funtor de Picard $\Pic_{X/k}$ de $X$ sobre $k$ es el
haz de grupos para la topología ($\fpqc$) de $k$ asociado al prehaz
\[
T\rightsquigarrow \Pic(X\times_kT)\qquad
(T\in\operatorname{Ob}(\Sch/k)).
\]
Cuando $k$ es algebraicamente cerrado, se sabe que el homomorfismo canónico
\[
\Pic(X)\longrightarrow\Pic_{X/k}(\Spec(k))
\]
es un isomorfismo.

% Macros locales de continuación
\providecommand{\Filt}{\operatorname{Fil}}
\providecommand{\fpqc}{\mathrm{fpqc}}
\providecommand{\Sch}{\operatorname{Sch}}
\providecommand{\Ob}{\operatorname{ob}}
\providecommand{\Gm}{\mathbb G_m}
\providecommand{\Der}{\operatorname{Der}}
\providecommand{\SL}{\operatorname{SL}}
\providecommand{\Aut}{\operatorname{Aut}}
\providecommand{\Mod}{\operatorname{Mod}}
\providecommand{\Tors}{\operatorname{Tors}}
\providecommand{\Mat}{\operatorname{Mat}}
\providecommand{\len}{\operatorname{long}}
\providecommand{\id}{\operatorname{id}}
\providecommand{\Spec}{\operatorname{Spec}}
\providecommand{\Picf}{\operatorname{Pic}}
\providecommand{\Ens}{\operatorname{Ens}}
\providecommand{\clp}{\operatorname{cl}_{P}}
\providecommand{\clo}{\operatorname{cl}^{\circ}}
\providecommand{\clb}{\operatorname{cl}_{\bullet}}
\providecommand{\Grb}{\operatorname{Gr}_{\bullet}}
\providecommand{\Gro}{\operatorname{Gr}^{\bullet}}
\providecommand{\longO}[1]{\operatorname{length}_{#1}}

\subsection*{2.4, continuación}

Cuando $X$ es propio sobre $k$, se sabe [4] que $\Pic^0_{X/k}$ es
representable por un esquema de grupos localmente de tipo finito
$\Pic_{X/k}$ sobre $k$. La componente neutra $\Pic^0_{X/k}$ de
$\Pic_{X/k}$ es un grupo algebraico conmutativo conexo sobre $k$.

\textbf{Definición 2.4.1.} Sea $X$ un esquema propio sobre un cuerpo $k$,
$L$ un $\mathcal O_X$-módulo invertible (resp. $x$ un elemento de
$\Pic(X)$). Se dice que $L$ (resp. $x$) es $\tau$-equivalente a cero si
existe un entero $n\ne0$ tal que $L^{\otimes n}$ (resp. $nx$) es
algebraicamente equivalente a cero. Los elementos de $\Pic(X)$
algebraicamente equivalentes (resp. $\tau$-equivalentes) a cero forman un
subgrupo $\Pic^0(X)$ (resp. $\Pic^\tau(X)$) de $\Pic(X)$.

\textbf{Lema 2.4.3.} Para todo esquema propio $X$ sobre un cuerpo
algebraicamente cerrado $k$, el grupo $\Pic^0(X)$ es divisible por todo
entero $n$ primo con la característica de $k$.

\emph{Demostración:} Como $k$ es algebraicamente cerrado, se tiene
$\Pic^0(X)\cong \underline{\Pic}^0_{X/k}(k)$. Puesto que $\Pic^0_{X/k}$ es conexo y de
tipo finito, el endomorfismo potencia $n$-ésima es sobreyectivo cuando $n$
es primo con la característica de $k$ (SGA 3 XV 1.3 y VI 1.3.2). Pero $k$
es algebraicamente cerrado; por tanto, el endomorfismo correspondiente del
grupo $\Pic^0_{X/k}(k)$ es sobreyectivo, C.Q.F.D.

De 2.4.3 se concluye que, en la sucesión exacta
\[
0\longrightarrow\Pic^0(X)\longrightarrow\Pic^\tau(X)
\longrightarrow T\longrightarrow0,
\]
$\Pic^0(X)$ es divisible por un número primo $p$ y $T$ es un grupo de
torsión, de modo que se puede aplicar 2.1.3 a la restricción de la aplicación
polinómica (2.3.1) a $\Pic^\tau(X)$. Se obtiene así:

\textbf{Proposición 2.4.4.} La aplicación polinómica $s_y$ (2.3.1) es
constante en las clases módulo $\Pic^\tau(X)$.

\emph{Demostración:} Si $K$ es una extensión algebraicamente cerrada de
$k$, se obtiene un diagrama conmutativo
\[
\begin{tikzpicture}[>=stealth,baseline=(current bounding box.center)]
\node (PicX) at (0,2) {\(\Pic(X)\)};
\node (Kk) at (3.7,2) {\(K_\bullet(\Spec(k))\)};
\node (Z) at (5.35,1) {\(\mathbb Z\)};
\node (Q) at (6.6,1) {\(\mathbb Q\)};
\node (PicXK) at (0,0) {\(\Pic(X_K)\)};
\node (KK) at (3.7,0) {\(K_\bullet(\Spec(K))\)};
\draw[->] (PicX) -- node[above] {\(s_y\)} (Kk);
\draw[->] (PicX) -- (PicXK);
\draw[->] (PicXK) -- node[below] {\(\underline s_y\)} (KK);
\draw[->] (Kk) -- (Z);
\draw[->] (KK) -- (Z);
\draw[right hook->] (Z) -- (Q);
\end{tikzpicture}
\]
donde $\underline s_y$ es la aplicación polinómica correspondiente a $s_y$ sobre
$\Pic(X_K)$. Se puede, por tanto, suponer que $k$ es algebraicamente
cerrado. La afirmación se deduce entonces inmediatamente de 2.1.3 y 2.4.3.

\textbf{Observación 2.4.5.} En la Exp. XIII se demostrará este resultado sin
utilizar la representabilidad del funtor $\Pic_{X/k}$.
